\chapter{Land and food}
\label{ch:land}

\section{The failure, stated exactly}

A tropical island with deep soils, adequate rainfall, a year-round growing
season, no frost, a trained agronomic profession and a hundred and fifty years of
commercial agriculture imports somewhere between two thirds and four fifths of
what it eats, cannot pay for the imports, and has between one and two million
hectares of farmland under \emph{marab\'u} scrub.\unv{}

This is the most complete policy failure in the book. It is more complete than
the energy failure, because energy at least has a physical explanation --- the
plants are old and the fuel is foreign. Nothing physical explains this one. Cuba's
neighbours with worse soil, worse rainfall and less education feed themselves.

It is also the clearest case where the embargo explains very little. Cuba may buy
agricultural inputs from anyone in the world, and does buy food from the United
States itself, for cash, under the exemption that has stood since 2000. The
constraint is not permission to import fertiliser. It is that Cuban farmers have
been given every incentive not to grow anything.

\section{Why}

Four mechanisms, in order of how much they explain.

\emph{Acopio.} The state procurement monopoly requires producers to sell a
defined share of output to the state at a state-set price, historically below the
cost of production and paid late, in a currency that was losing value while the
farmer waited. A farmer facing that has three rational responses: produce the
minimum the quota requires, divert everything else to the informal market, or
stop. All three are observed, and all three are then prosecuted as indiscipline.

\emph{Tenure.} The usufruct grants of 2008 and 2012 handed out large areas of
idle land, and they handed out the weakest possible interest in it. The land
cannot be sold. It cannot be mortgaged. Its inheritance is conditional. The term
is limited and renewal is discretionary. Buildings erected on it may belong to
the state. Nobody plants a fruit tree that yields in year seven on a tenure that
might end in year five, and nobody can borrow against land a bank cannot take.
The result is that the grants produced short-cycle crops on the best parcels and
scrub on the rest.

\emph{Inputs and equipment.} No functioning market in fertiliser, seed,
veterinary supplies, fencing, irrigation pipe, tyres or spares. What exists is
allocated administratively, which means it is allocated to state farms and to
whoever knows the person allocating.

\emph{Everything after the harvest.} No transport, no refrigeration, no
processing, no packaging. Cuban post-harvest losses are very large --- estimates
of a fifth to a third of what is grown, and higher for perishables --- and the
loss happens between a field where the food exists and a city where people are
hungry.\unv{} A cold chain is not agriculture policy in the usual sense; in Cuba
it is the largest single available increase in food supply, and it requires
electricity, which is why this chapter follows Chapter~\ref{ch:energy}.

Every one of these four is domestic, legal and reversible.

\section{The design}

\begin{design}{Tenure}
\label{d:tenure}
\begin{rules}
\rl{1}{Convert usufruct to ownership} Existing usufruct holders who have worked
their land for a qualifying period receive full, registered, transferable
freehold or a long lease of not less than ninety-nine years, at nominal
consideration. Not a longer usufruct --- an interest that can be sold, mortgaged,
inherited and defended in court.

\rl{2}{Registered} Every parcel entered on the public cadastre and registry of
Chapter~\ref{ch:capital}, with boundaries surveyed and title guaranteed by the
state. An unregistered right is a right that has to be re-argued with each
official.

\rl{3}{Mortgageable} Agricultural land may secure a loan and may be taken by a
creditor on default. This is the clause that makes rural credit possible, it is
the clause that will be resisted hardest, and the fear behind the resistance ---
that farmers lose their land --- is met by clause~4 rather than by refusing to
allow borrowing at all.

\rl{4}{Anti-concentration} A published ceiling on holdings by a single beneficial
owner, and a right of first refusal for the sitting occupier and for neighbouring
smallholders on any sale. Cuba's land question was settled in 1959 against
latifundia and this design does not reopen it; the objective is a market in
family-scale and cooperative holdings, not the reassembly of the sugar estates.

\rl{5}{Foreign ownership} Foreigners may lease agricultural land, and may not
own it. Chapter~\ref{ch:capital} takes a permissive line on foreign investment
generally, and this is the deliberate exception, for reasons that are political
rather than economic and that any Cuban government would face.

\rl{6}{Use it} Land held out of production beyond a defined period is subject to
a rising land tax under Design~\ref{d:tax}.4 rather than to confiscation. A tax
is administrable, is not discretionary, and does not require an official to
decide whether a farm is being farmed.
\end{rules}
\end{design}

\begin{design}{Markets}
\label{d:acopio}
\begin{rules}
\rl{1}{End acopio} The state procurement monopoly is abolished. Producers sell to
whom they choose at prices they agree. The state continues to buy --- for
schools, hospitals, the armed forces and the strategic reserve --- as one buyer
among others, at market prices, paying on published terms and on time.

\rl{2}{Input imports opened} Any producer, cooperative or trading firm may import
fertiliser, seed, veterinary products, machinery, spares and irrigation
equipment, free of licence and of duty, per Design~\ref{d:commit}.2.

\rl{3}{Wholesale markets} Physical wholesale markets in each province, open to
any buyer and seller, with published prices. The absence of a wholesale tier is
why a Cuban restaurant buys its vegetables in a retail shop in competition with
households, and why the private sector gets blamed for shortages it is standing
in the queue for.

\rl{4}{Rural credit} Agricultural lending against the registered, mortgageable
title of Design~\ref{d:tenure}, through banks and credit cooperatives, with a
partial public guarantee in the first phase to establish the instrument. Cuba
currently has neither the collateral nor the lenders.

\rl{5}{No price controls} Retail price caps on food are prohibited. They were
imposed on the private importers from 2024, they produced exactly the withdrawal
of supply that price caps always produce, and the distributional problem they
were addressing is met on the spending side by Design~\ref{d:libreta} instead.
\end{rules}
\end{design}

\section{Food security is not self-sufficiency}

A distinction the Cuban policy tradition has consistently confused, at
considerable cost.

Self-sufficiency means producing domestically what one consumes. Food security
means that people can reliably obtain food. These are different objectives with
different instruments, and pursuing the first is a poor way to achieve the
second on a small island.

Cuba should not try to grow its own wheat --- it cannot, at any price. It should
not try to compete in bulk maize or soy against continental producers with a
tenth of its transport cost per tonne. Those are imports, permanently, and the
security question about them is whether the country can pay, which is a question
about export earnings and reserves rather than about agronomy.

What Cuba can produce competitively is a well-defined list: roots and tubers,
plantain, vegetables and salads, tropical fruit, beans, eggs, poultry, pork,
dairy at a modest scale, fish and aquaculture, and the high-value export crops it
has always had --- tobacco, coffee, cacao, honey, citrus, rum. Several of these
are worth far more per hectare than anything the country currently grows, and
Cuban coffee and cacao in particular sell at premiums in specialty markets that
Cuba has never organised itself to supply.

\begin{design}{What the state is actually for}
\label{d:foodsec}
\begin{rules}
\rl{1}{Reserves, not targets} Food security is delivered by a strategic reserve
of staples sized in months of consumption, published, rotated, and financed from
the stabilisation fund --- not by production targets for crops the country grows
badly.

\rl{2}{Buy where it is cheap} Import policy is set on cost, not on the political
identity of the seller. Cuba currently buys food from the United States for cash
while describing it as a hostile power, which is awkward and correct.

\rl{3}{Public goods} The state's agricultural role is research and extension,
plant and animal health and quarantine, cadastre and registry, rural roads,
irrigation infrastructure, and market price information. These are the things
private actors will not provide and the things the Cuban state has stopped doing
while it was busy setting prices.

\rl{4}{Export promotion} Certification, traceability, phytosanitary compliance
and origin protection for the export crops --- which is what standing between a
Cuban cacao farmer and a European specialty buyer, and what no individual farmer
can do alone.
\end{rules}
\end{design}

\section{Sugar}

The honest answer is that it does not come back, and a design that plans on its
return is planning on nostalgia.

Cuban sugar was eight and a half million tonnes in 1970 and is now measured in
the low hundreds of thousands. The mills are gone or cannibalised, the railways
that served them are gone, the cane fields are under scrub, the workforce has
aged out, and the world price is set by Brazilian and Indian producers operating
at a scale and cost Cuba cannot approach. Rebuilding the industry would require
capital that has better uses in every other chapter of this book.

What is worth keeping is smaller and more specific. The remaining efficient mills
in the best cane districts, run as an integrated sugar-ethanol-bagasse-power
business rather than as raw sugar exporters, on a scale of perhaps a fifth of the
1989 industry. The rum, which is a high-value branded export with a genuine
international market and the one part of the sugar complex that has never
stopped earning. And the land, which is the real asset: a million hectares of
former cane land, flat, road-served and now idle, is exactly what
Chapter~\ref{ch:energy} needs for photovoltaics and what Design~\ref{d:tenure}
needs for redistribution.

The social question is the difficult part and it should not be waved away. The
sugar towns --- and there are hundreds of them, built around a single mill, with
no other employer --- were devastated by the closures of 2002 and have never
recovered. Whatever is decided about the industry, those communities need an
answer, and the honest one is that it is not sugar. It is the eastern weighting
of Design~\ref{d:energyseq}.3 and Design~\ref{d:dutch}.4, the equalisation formula
of Design~\ref{d:municipal}.2, and land.

\section{What the 1990s did and did not prove}

The organic and urban agriculture that emerged from the Special Period is the
most internationally celebrated thing in Cuban agriculture and the most
over-claimed.

What it demonstrably did: put fresh vegetables into cities at a moment when the
transport system had failed, using raised beds on vacant urban lots, biological
pest control and composting; and it did so on a scale that supplied a meaningful
share of Havana's fresh vegetable consumption. That is a real achievement, it was
achieved by Cuban agronomists under extreme constraint, and the technique
transfers.

What it did not do: feed the country. Cuba's food import dependence rose across
the whole period rather than falling. Urban horticulture is good at leafy
vegetables and herbs on small plots close to consumers, and it is not a
substitute for field crops, protein or calories. And the low-input methods were
adopted because there were no inputs, not because they outperformed; where
fertiliser has been available since, farmers have used it.

The design position is therefore neither the romantic one nor the dismissive one.
Keep the urban horticulture, which works and is cheap, and extend it under
Design~\ref{d:municipal}. Keep the biological pest control and the agronomic
extension capacity, which are genuinely good. And do not build national food
policy on a legend, because Cubans have been hungry throughout the entire period
during which the legend has been circulating abroad.
